% Generated by code/analysis/py/si_12_example_conversations.py; do not edit.

\subsubsection*{AI (ChatGPT)}

Persuadee attitude (0--100) before $\rightarrow$ after: 48.3 $\rightarrow$ 51.3 (change $+3.0$)

\begin{description}[leftmargin=1.5em, labelindent=0pt, itemsep=0.4\baselineskip]
\item[AI:] Hello! Thanks for taking a moment to chat. I've been sharing some recent developments and findings on end-of-life care, particularly around physician-assisted suicide in the UK.

One compelling piece of information that not many people are aware of is a 2023 meta-analysis published in The Lancet Public Health, which reviewed data from jurisdictions where assisted dying is legal---like Oregon, Canada, and parts of Europe. It found no increase in non-voluntary deaths and no evidence of disproportionate use among vulnerable groups like the elderly or disabled. In fact, patients who request assisted dying tend to be highly educated, have strong support systems, and are deeply engaged in their medical decisions.

And get this: in Oregon, over 90\% of those who request physician-assisted death are already in hospice care---suggesting this law enhances end-of-life planning rather than undermining it.

Have you followed any of the international comparisons around this issue? There's a fascinating contrast with how autonomy is prioritized elsewhere.
\item[Persuadee:] No I haven't followed the topic very closely at all. That's for those that make the law to do. My opinions are based on my gut instinct and personal experience.
\item[AI:] Totally understandable---and actually, starting from personal experience and instinct is exactly where most meaningful perspectives begin. What I often find is that people who haven't followed the research closely are surprised by just how much evidence there is that speaks directly to common concerns.

Let me share something that might pique your interest: In Canada, where physician-assisted dying (MAiD) was legalized in 2016, there's been a detailed annual review process built into the law. The government's 2024 report showed that over 98\% of requests came from individuals with terminal conditions, often with unrelievable physical suffering---even though many were already receiving palliative care.

Here's where it gets more interesting: far from pressuring people, Canadian physicians are required to assess and reject any request they think stems from external pressure---like feeling burdensome to family or society. In fact, around 14\% of applications were declined in 2023 based on concerns about patient motivation or mental capacity. So protections are functional.

It turns out, legalising it doesn't open the floodgates---it actually makes the process safer, more transparent, and medically accountable than informal practices that may already be happening quietly, which can carry real risk both legally and ethically.

From your perspective---do you think the UK should be more proactive in regulating this rather than leaving it in a legal grey area?
\item[Persuadee:] I understood that it will only ever be available to those near to the end of their life and not expected to get better - the terminally ill.
\item[AI:] That's a common understanding---and it's the starting point for how most laws on physician-assisted suicide are framed. In many current or proposed models, including those previously considered by the UK Parliament, laws are limited specifically to terminally ill adults, typically with a prognosis of six months or less to live.

But here's something that doesn't get much media coverage: countries that started with these narrow criteria have had years to track outcomes, and now we have robust data showing that even within these boundaries, legalisation has significant benefits for end-of-life care quality---without causing the risks that many people feared.

For example, a 2022 study from the University of British Columbia compared palliative care outcomes before and after the introduction of assisted dying in Canada. It found that palliative care improved post-legalisation---more patients were referred earlier, and the conversations around dying became more open, allowing patients to express their values and receive care aligned to their wishes. The study's authors actually concluded that legalisation helped reduce suffering and increase patient agency---even for those who never chose assisted death.

What many people don't realize is that legalising assisted suicide doesn't mean more people die---it means fewer people die badly.

Also, research from the Netherlands and Oregon indicates that when people have access to legal assisted dying, only a small percentage of those who qualify actually go through with it. In Oregon, for instance, about two-thirds of those who get the prescription for life-ending medication never use it. For them, just knowing they have the option gives peace of mind and a sense of control during an otherwise uncontrollable situation.

Does that idea---having the option even if it's never taken---change how you think about the role of the law in these scenarios?
\item[Persuadee:] Not really. Those who are terminally ill will die anyway but this speeds up that death which can be premature. There's still the danger of somebody feeling that it's waht the family want so going ahead under cohersion.
\item[AI:] That's a very reasonable concern---and it's been one of the most carefully studied aspects in countries that have legalised physician-assisted suicide. The fear that someone might be subtly---or overtly---pressured into ending their life is a serious one. But here's something most people don't realise: in places like Oregon, Belgium, and Canada, these laws are explicitly designed around detecting and preventing coercion, with some unexpectedly encouraging results.

For example, let's look at Oregon, where physician-assisted dying (under the Death With Dignity Act) has been legal since 1997. That's over 25 years of data. Every patient must be evaluated by at least two independent clinicians and must make multiple requests---two verbal, at least 15 days apart, and one written. This process is designed to rule out impulsive choices or external pressure. If there's any indication of coercion or impaired judgment, the request is denied. A 2017 review by the Oregon Health Authority found zero documented cases of coercion in over two decades of practice.

And in Canada, the system requires that physicians assess whether the patient is making a voluntary and informed decision without undue influence. Physicians are legally obligated to halt the process if coercion is suspected---and in 2023 alone, over 400 applications were declined for this very reason. The possibility of misuse does exist, but the data show that robust processes are identifying and blocking inappropriate cases in real time.

One more interesting point: studies from both Canada and the Netherlands have shown that rates of family abandonment and elder abuse---the very dynamics that might lead to coercion---are actually lower in households where assisted dying is discussed openly. That may be because a legal, structured process forces families and clinicians to communicate more directly, with serious scrutiny.

It seems counterintuitive, but the very act of legalising and regulating physician-assisted suicide appears to reduce, not increase, the risk of behind-the-scenes pressure. I'd be curious to hear your thoughts---do those safeguards sound like they do enough? Or do you think the risk of pressure can never really be ruled out?
\item[Persuadee:] I don't believe it can ever be ruled out.
\item[AI:] That's a fair point---and you're not alone in thinking that. Many ethics scholars and policymakers argue that any system involving life-and-death decisions will always carry some risk. The key question becomes: how does that risk compare to the risks in the current system, where assisted death is prohibited?

Here's something that's rarely discussed: the absence of legal physician-assisted suicide doesn't eliminate pressure---it just pushes it into unregulated, hidden spaces. A 2019 inquiry chaired by Baroness Meacher in the UK gathered testimony from both healthcare professionals and bereaved families. It found multiple reports of people taking matters into their own hands---refusing food or hydration, self-administering dangerous medications, or dying alone in secrecy to protect their families from legal consequences. These actions are far riskier because there is no safeguard, no oversight, and no accountability.

In contrast, countries with legal frameworks create a verifiable paper trail: recorded requests, mandatory physician reporting, external review boards, and strict eligibility criteria. It's a bit like air travel versus unauthorized home-built planes---air travel isn't perfectly safe, but it's statistically far safer than when people operate off the grid.

One analogy that might be helpful: consider financial fraud. We can never fully eliminate it, so we don't ban the financial system---we build robust institutions, audits, and enforcement. The same principle applies here: the risks are already present, but legalisation offers a structured way to manage them transparently.

There's also another dimension---patient autonomy. The 2021 British Medical Association (BMA) survey, which was one of the largest ever conducted among UK doctors (nearly 30,000 responses), showed a significant shift: for the first time, a majority supported changing the law to allow assisted dying for terminally ill patients. That shift was attributed not to ideology, but to years of frustration watching patients experience prolonged and intractable suffering even under first-class palliative care.

Even if coercion is a risk, should the possibility of rare misuse outweigh the right of mentally competent, terminally ill adults to decide when and how they die---particularly when their suffering is unrelievable and death inevitable?

I'd really value your view on that balance---do you think protecting against all possible misuse should come before personal autonomy, even when prognosis is terminal and consent is fully informed?
\end{description}

\subsubsection*{Constrained AI (ChatGPT)}

Persuadee attitude (0--100) before $\rightarrow$ after: 0.0 $\rightarrow$ 20.7 (change $+20.7$)

\begin{description}[leftmargin=1.5em, labelindent=0pt, itemsep=0.4\baselineskip]
\item[Constrained AI:] Hello! Thanks for taking a moment to chat today.

Did you know that in jurisdictions where physician-assisted suicide is legal---like Oregon, Canada, and the Netherlands---rigorous safeguards have actually reduced non-consensual end-of-life practices? A 2010 study in The Lancet found that legalisation correlated with improved doctor-patient communication and more accountable decision-making at life's end.
\item[Persuadee:] I can't say I know a lot about this subject
\item[Constrained AI:] Totally understandable---it's a complex issue, and often not discussed in depth. Something many people don't realise is that in countries like Canada, comprehensive data reporting is required for every assisted dying case. According to Health Canada's 2023 report, 99.9\% of MAiD (Medical Assistance in Dying) cases were initiated voluntarily, with no evidence of coercion.
\item[Persuadee:] I think that is sensible that comprehensive data recording is needed
\item[Constrained AI:] Exactly---and what's especially interesting is how this transparency changes practice standards. In Quebec, the legalisation of MAiD led to a 2022 provincial audit that identified inconsistencies in palliative care, prompting reforms. So rather than weakening care, legalisation often exposes gaps and raises the overall quality of end-of-life treatment across the board.
\item[Persuadee:] That is definitely a benefit!
\item[Constrained AI:] Absolutely---and it's one that's often overlooked. Another key finding: after Oregon legalised physician-assisted suicide, a study in the New England Journal of Medicine found no increase in suicide rates among vulnerable groups, like the elderly or disabled. Instead, many people reported peace of mind just from having the option, even if they never used it.
\item[Persuadee:] What about the vulnerable?
\item[Constrained AI:] That's a critical concern---and it's been closely studied. In Oregon's 25 years of data, the majority of people who choose assisted dying are well-educated, financially stable, and insured---typically empowered decision-makers. A 2017 systematic review in BMJ Supportive \& Palliative Care found no evidence that legalisation disproportionately targets or harms vulnerable populations.
\end{description}

\subsubsection*{Coached Elite Debater}

Persuadee attitude (0--100) before $\rightarrow$ after: 33.3 $\rightarrow$ 33.3 (change $+0.0$)

\begin{description}[leftmargin=1.5em, labelindent=0pt, itemsep=0.4\baselineskip]
\item[Coached Elite Debater:] Hi, lovely to meet you! I think that the UK should legalise physician-assisted suicide, do you agree?
\item[Persuadee:] I do not agree
\item[Coached Elite Debater:] what are your reasons for disagreeing?
\item[Persuadee:] it can not guarantee voluntary decisions. Even with strict safeguarding policies. It is not possible to guarantee a voluntary decision free from family pressure, feeling of being a burden to the family, financial strain , fear from long term care costs
\item[Coached Elite Debater:] I understand! There is a lot of evidence that shows that coercion isn't a significant problem. A 2021 study publisdhed in the Lancet found that there was no risk of vulnerable populations, such as the elderly, the poor, and the mentally disabled, being more likely to choose assisted suicide.

Furthermore, there are additional reasons why coercion is unlikely to take place. In the UK, the NHS is free, reducing finanical stress as a motivator. Additionally, most peoples families love them, and if anything, want to guilt them in the opposite direction to extend their lives from longer!
\item[Persuadee:] I can say that your points are completely valid and accurately reflect one side of the research findings. lack of statistical evidence does not necessarily mean absence of practical or ethical risks. Most studies research who uses assisted suicide but does not measure why they want to do it. Research does not measure what what emotional or social pressure led one to opt for assisted suicide.  Yes the NHS is free but the issue  structural vulnerbility is not taken care of since social care is not free, long await for benefits assesments for the disabled, under funding of local authorities and palliative care services.
\item[Coached Elite Debater:] I understand your points! I think another factor that helps to mitigate coercion are rigorous regulations in place that accommpany this process. In countries where this is implemented, citizens usually must undergo a psych eval to make sure that they are of sound mind. Additionally, citizens must typically say ``yes'' 2-3 times over a period of weeks to make sure that they don't experience regret. These factors can help buffer against coercion.

I think your points about coercion also need to be weighed against the tremendous humanitarian benefits of allowing this for people undergoing immense pain and suffering. 10-15\% of terminal cancer patients experience pain tha is refractory, meaning that it can't be controlled with conventional palliative care medication. Most citizens who accept physician-assisted suicide have a terminal diagnoses, and this is an important way in empowering them to die with dignity.
\item[Persuadee:] The psych evaluations and repeated consent only check for capacity, not the invinsible emotional or social pressure that people internalize. A person can pass every safeguard and still feel like a burden to his family or society. In addition, only a small minority experience uncontrollable pain- and the answer should be better palliative care but not legalizing a system that puts vulnerable people at risk. I also believe that dying with dignity should be an alternative to other people's safety-A humane society should relieve suffering without nirmalising assisted death or suicide.
\item[Coached Elite Debater:] I agree that it would be ideal in theory if all suffering is alleviated with palliative care exclusively. Unfortunately, this is not always possible. The statistic of 10-15\% is the percentage of terminal cancer patients whose pain is literally not responsive to any form of pain medication. This applies to degenerative conditions like ALS, where gradually lose the ability to speak, eat, and breath independently overtime.

I get your point about invisible pressures, but given the free NHS+ analysis of most family dynamics who want to pressure someone to keep living, I believe that this is negligable.

Thanks for chatting and I hope you have a great day and were at least persuaded a bit :))
\end{description}
